\chapter{Space, alternation, and QBF}
\label{chap:space}

This chapter covers space-bounded computation: the parametric classes
$\cc{DSPACE}$ and $\cc{NSPACE}$, the logarithmic-space classes $\cc{L}$,
$\cc{NL}$, $\cc{coNL}$ and $\cc{FL}$, polynomial space, and quantified Boolean
formulas. Space is measured by an honest auxiliary-space convention in which
neither the input tape nor the output tape can serve as hidden workspace.
Formalized so far: the elementary containments, closure of $\cc{PSPACE}$ under
complement, regular languages in $\cc{L}$, $\cc{L} \subseteq \cc{P}$,
$\cc{FL} \subseteq \cc{FP}$, $\cc{NL} \subseteq \cc{P}$, Savitch's theorem in
the polynomial-union form $\cc{PSPACE} = \cc{NPSPACE}$,
$\cc{PSPACE} \subseteq \cc{EXP}$, the syntax and semantics of quantified
Boolean formulas, space-aware Hoare contracts, and a reduction of the
Immerman--Szelepcs\'enyi theorem to the construction of one counting machine.
Savitch's theorem and $\cc{IP} \subseteq \cc{PSPACE}$ both go through one
lever: iterating a polynomial-time step function on a polynomially bounded
state stays in $\cc{PSPACE}$. The main open directions are the unconditional
$\cc{NL} = \cc{coNL}$, a log-space programming layer, PSPACE-completeness of
TQBF, alternation, and the space hierarchy theorems. Throughout, $\log n$ is
$\lfloor \log_2 n \rfloor$ and $f = O(g)$ is eventual domination up to a
constant factor.

\section{Space-bounded deciders and classes}

\begin{definition}[Space-bounded deciders and transducers]
  \label{def:space-decides-in-space}
  \lean{Complexity.Cfg.WithinAuxSpace, Complexity.Cfg.WithinDecisionSpace,
    Complexity.TM.DecidesInSpace, Complexity.NTM.DecidesInSpace,
    Complexity.TM.IsTransducer, Complexity.NTM.IsTransducer,
    Complexity.TM.ComputesInSpace}
  \leanok
  \uses{def:tm, def:ntm, def:language}
  A configuration on an input of length $n$ is within auxiliary space $s$ when
  every work head is at position at most $s$ and the input head is at position
  at most $n + s + 1$: the input and its first trailing blank are free, and
  travel farther into the blank tail is charged. The output head is not
  constrained. A configuration is within decision space $s$ when in addition
  its output head is at position at most $s + 1$; only the verdict cell $1$ is
  free.

  A DTM decides $L$ in space $S$ if every configuration reachable from its
  initial configuration on $x$ is within decision space $S(|x|)$ and, for every
  $x$, some reachable configuration is halted with verdict $1$ if $x \in L$
  and $0$ if $x \notin L$. An NTM decides $L$ in space $S$ if there is a time
  bound $T$ such that every choice sequence halts within $T(|x|)$ steps,
  $x \in L$ exactly when some choice sequence of length $T(|x|)$ ends halted
  with verdict $1$, and every configuration along every path, up to time
  $T(|x|)$, is within decision space $S(|x|)$. A machine is a
  \emph{transducer} if its transition function never moves the output head
  left (for an NTM, under either choice bit); the two deciding predicates do
  not require this. A DTM computes $f$ in space $S$ if it is a transducer,
  every reachable configuration is within auxiliary space $S(|x|)$, and on
  every input $x$ it halts with output $f(x)$; the output length is not
  charged.
\end{definition}

\begin{definition}[Deterministic space]
  \label{def:dspace}
  \lean{Complexity.DSPACE}
  \leanok
  \uses{def:space-decides-in-space}
  $\cc{DSPACE}(S)$ is the class of languages decided by some DTM, with any
  number of work tapes, in space $f$ for some $f = O(S)$.
\end{definition}

\begin{definition}[Nondeterministic space]
  \label{def:nspace}
  \lean{Complexity.NSPACE}
  \leanok
  \uses{def:space-decides-in-space}
  $\cc{NSPACE}(S)$ is the class of languages decided by some NTM, with any
  number of work tapes, in space $f$ for some $f = O(S)$.
\end{definition}

\begin{definition}[Polynomial space]
  \label{def:PSPACE}
  \lean{Complexity.PSPACE}
  \leanok
  \uses{def:dspace}
  $\cc{PSPACE} = \bigcup_{k} \cc{DSPACE}(n^k)$.
\end{definition}

\begin{definition}[Nondeterministic polynomial space]
  \label{def:NPSPACE}
  \lean{Complexity.NPSPACE}
  \leanok
  \uses{def:nspace}
  $\cc{NPSPACE} = \bigcup_{k} \cc{NSPACE}(n^k)$.
\end{definition}

\begin{definition}[Logarithmic space]
  \label{def:L}
  \lean{Complexity.L}
  \leanok
  \uses{def:space-decides-in-space}
  $\cc{L}$ is the class of languages decided by a deterministic transducer,
  with any number of work tapes, in space $f$ for some $f = O(\log n)$.
\end{definition}

\begin{definition}[Nondeterministic logarithmic space]
  \label{def:NL}
  \lean{Complexity.NL}
  \leanok
  \uses{def:space-decides-in-space}
  $\cc{NL}$ is the class of languages decided by a nondeterministic
  transducer, with any number of work tapes, in space $f$ for some
  $f = O(\log n)$.
\end{definition}

\begin{definition}[Complement of NL]
  \label{def:coNL}
  \lean{Complexity.coNL}
  \leanok
  \uses{def:NL}
  $\cc{coNL}$ is the class of languages whose complements are in $\cc{NL}$.
\end{definition}

\begin{definition}[Log-space functions and search problems]
  \label{def:FL}
  \lean{Complexity.FL, Complexity.FNL, Complexity.TFNL}
  \leanok
  \uses{def:space-decides-in-space, def:L, def:pairlang}
  $\cc{FL}$ is the class of functions $\{0,1\}^* \to \{0,1\}^*$ computed by a
  DTM in space $S$ for some $S = O(\log n)$. The output is unbounded in length
  and one-way. $\cc{FNL}$ is the class of polynomially balanced relations whose
  pair language is in $\cc{L}$, and $\cc{TFNL}$ consists of the relations in
  $\cc{FNL}$ that are total: every $x$ has some $y$ with $R(x, y)$.
\end{definition}

\begin{definition}[Simultaneous time and space; SC]
  \label{def:space-dtisp}
  \lean{Complexity.TM.DecidesInTimeSpace, Complexity.DTISP, Complexity.SC}
  \leanok
  \uses{def:space-decides-in-space}
  A DTM decides $L$ in time $T$ and space $S$ if every reachable configuration
  on $x$ is within decision space $S(|x|)$ and on every input it halts within
  $T(|x|)$ steps with the correct verdict. $\cc{DTISP}(T, S)$ is the class of
  languages decided by a \emph{single} DTM in time $t$ and space $s$ for some
  $t = O(T)$ and $s = O(S)$. Steve's class is
  $\cc{SC} = \bigcup_{k, j} \cc{DTISP}(n^k, (\log n)^j)$.
\end{definition}

\section{Elementary containments and closure}

\begin{lemma}[Determinism embeds in nondeterminism]
  \label{lem:space-dspace-nspace}
  \lean{Complexity.TM.toNTM_decidesInSpace, Complexity.TM.toNTM_isTransducer,
    Complexity.DSPACE_subset_NSPACE, Complexity.DSPACE_mono,
    Complexity.NSPACE_mono, Complexity.L_subset_NL,
    Complexity.PSPACE_subset_NPSPACE}
  \leanok
  \uses{def:space-decides-in-space, def:dspace, def:nspace, def:L, def:NL,
    def:PSPACE, def:NPSPACE}
  If a DTM decides $L$ in space $S$, then so does its embedding as an NTM whose
  two transition functions coincide, and the embedding of a transducer is a
  transducer. Hence for every $S$, $\cc{DSPACE}(S) \subseteq \cc{NSPACE}(S)$.
  If $S_1 = O(S_2)$ then $\cc{DSPACE}(S_1) \subseteq \cc{DSPACE}(S_2)$ and
  $\cc{NSPACE}(S_1) \subseteq \cc{NSPACE}(S_2)$. Consequently
  $\cc{L} \subseteq \cc{NL}$ and $\cc{PSPACE} \subseteq \cc{NPSPACE}$.
\end{lemma}
\begin{proof}
  \leanok
  The embedded NTM follows the same run on every choice sequence; its uniform
  time bound is the largest halting time over the finitely many inputs of each
  length. Monotonicity is transitivity of $O$.
\end{proof}

\begin{proposition}[Time bounds space]
  \label{prop:space-dtime-dspace}
  \lean{Complexity.DTIME_subset_DSPACE, Complexity.DTIME_subset_NSPACE,
    Complexity.P_subset_PSPACE, Complexity.P_subset_NPSPACE}
  \leanok
  \uses{def:dtime, def:dspace, def:nspace, def:P, def:PSPACE, def:NPSPACE}
  For every $T$, $\cc{DTIME}(T) \subseteq \cc{DSPACE}(T)$ and
  $\cc{DTIME}(T) \subseteq \cc{NSPACE}(T)$. Hence
  $\cc{P} \subseteq \cc{PSPACE}$ and $\cc{P} \subseteq \cc{NPSPACE}$.
\end{proposition}
\begin{proof}
  \leanok
  \uses{lem:space-dspace-nspace}
  Every head moves at most one cell per step and a decider in time $f$ halts
  within $f(|x|)$ steps, so every reachable configuration of the same machine
  is within decision space $f(|x|)$. The nondeterministic statements compose
  this with Lemma~\ref{lem:space-dspace-nspace}.
\end{proof}

\begin{lemma}[Closure under complement]
  \label{lem:space-compl}
  \lean{Complexity.DSPACE_compl, Complexity.PSPACE_compl}
  \leanok
  \uses{def:dspace, def:PSPACE}
  If the constant function $1$ is $O(S)$, then $\cc{DSPACE}(S)$ is closed under
  complement. In particular $\cc{PSPACE}$ is closed under complement.
\end{lemma}
\begin{proof}
  \leanok
  Run the decider, rewind the output head to the verdict cell, and flip the
  bit. The rewind costs one extra cell, so the new machine decides the
  complement in space $f + 1$; the delicate point is that the space predicate
  constrains every reachable configuration, not only the last.
\end{proof}

\begin{lemma}[PSPACE as an explicit polynomial window]
  \label{lem:space-poly-window}
  \lean{Complexity.exists_poly_window_of_mem_PSPACE,
    Complexity.mem_PSPACE_of_polyWindow}
  \leanok
  \uses{def:PSPACE}
  $L \in \cc{PSPACE}$ if and only if there are a DTM $M$ and a polynomial $q$
  with natural-number coefficients such that every configuration $M$ reaches
  on $x$ is within decision space $q(|x|)$ and $M$ halts on every input with
  the correct verdict.
\end{lemma}
\begin{proof}
  \leanok
  Unfold the union and replace the asymptotic bound by a dominating polynomial.
  The iteration lemma, Savitch's theorem, and the
  $\cc{PP} \subseteq \cc{PSPACE}$ and $\cc{PH} \subseteq \cc{PSPACE}$ machines
  enter $\cc{PSPACE}$ through the converse direction.
\end{proof}

\begin{theorem}[Regular languages are in $\cc{L}$]
  \label{thm:space-regular-in-L}
  \lean{Complexity.mem_DSPACE_zero_of_isRegular,
    Complexity.mem_DSPACE_of_isRegular, Complexity.mem_L_of_isRegular,
    Complexity.mem_L_of_nfa, Complexity.mem_L_of_twoWayNA,
    Complexity.mem_L_of_finite_monoid, Complexity.mem_L_of_regex}
  \leanok
  \uses{def:dspace, def:L}
  Every regular binary language, in the sense shared by Mathlib and CSLib
  (accepted by a finite deterministic automaton), is in $\cc{DSPACE}(0)$, hence
  in $\cc{DSPACE}(S)$ for every $S$, and is in $\cc{L}$. The languages of
  CSLib's finite nondeterministic automata, one-way or two-way, the languages
  of regular expressions over $\{0,1\}$, and the preimages of subsets of
  finite monoids under monoid homomorphisms from binary strings are in
  $\cc{L}$.
\end{theorem}
\begin{proof}
  \leanok
  The finite-state scanner of Theorem~\ref{thm:polytime-regular-in-P} has no
  work tapes, never moves its input head past the first blank after the input,
  and never moves its output head past the verdict cell, so it decides the
  language in space $0$. Its output head never moves left, so it is also a
  transducer, as membership in $\cc{L}$ requires. For the other models, CSLib
  first shows that the language is regular.
\end{proof}

\section{Configuration graphs and space-to-time}

\begin{definition}[Configuration graph]
  \label{def:space-config-graph}
  \lean{Complexity.NTM.Succ, Complexity.NTM.ReachesCfg,
    Complexity.NTM.ReachesCfgIn, Complexity.NTM.ReachesCfgLe,
    Complexity.NTM.reachSet}
  \leanok
  \uses{def:ntm}
  The configuration graph of an NTM has an edge from each non-halted
  configuration to each of its two successors, one per choice bit.
  $\mathrm{ReachesCfg}$ is the reflexive-transitive closure of this relation,
  $\mathrm{ReachesCfgIn}(t)$ and $\mathrm{ReachesCfgLe}(t)$ are reachability in
  exactly and in at most $t$ steps, and $\mathrm{reachSet}(c_0, t)$ is the set
  of configurations reached from $c_0$ within $t$ rounds of successor closure,
  that is, the state of breadth-first search after $t$ rounds.
\end{definition}

\begin{definition}[Windowed configurations and their codes]
  \label{def:space-bounded-configs}
  \lean{Complexity.Windowed, Complexity.windowed_init, Complexity.Code,
    Complexity.cfgCode, Complexity.cfgCode_inj, Complexity.card_Code,
    Complexity.card_Code_le_two_pow, Complexity.Windowed.step,
    Complexity.Windowed.stepCfg, Complexity.windowed_of_reachesCfg}
  \leanok
  \uses{def:space-decides-in-space, def:space-config-graph}
  Fix an input $x$ of length $n$ and a space bound $s$, and consider
  configurations with state set $Q$ and $k$ work tapes. A configuration is
  \emph{windowed} if its input tape holds the initial input contents, every
  work cell beyond position $s$ is blank, and every output cell beyond
  position $s + 1$ is blank; the initial configuration is windowed. The code
  of a configuration records its state, its input head clamped to at most
  $n + s + 1$, each work head clamped to at most $s$ together with work cells
  $0, \dots, s$, and its output head clamped to at most $s + 1$ together with
  output cells $0, \dots, s + 1$. Two windowed configurations within decision
  space $s$ with the same code are equal. There are exactly
  $|Q|\,(n + s + 2)\,\bigl((s + 1)\,4^{s + 1}\bigr)^k\,(s + 2)\,4^{s + 2}$
  codes, which is at most $2^{|Q| + (n + s + 2) + 3k(s + 1) + 3(s + 2)}$. A DTM
  step, or an NTM step under either choice bit, from a windowed configuration
  within decision space $s$ gives a windowed configuration. Hence, if every
  configuration reachable in the configuration graph from a windowed
  configuration is within decision space $s$, then all of them are windowed.
  These declarations live in proof-internal modules and are shared by
  $\cc{PSPACE} \subseteq \cc{EXP}$, $\cc{NL} \subseteq \cc{P}$, Savitch's
  theorem and the inductive-counting certificate. The sharper count for
  transducers under auxiliary space, in Theorem~\ref{thm:space-space-to-time},
  is separate.
\end{definition}

\begin{theorem}[Space bounds give time bounds]
  \label{thm:space-space-to-time}
  \lean{Complexity.TM.transducerConfigBound,
    Complexity.TM.card_transducerSnapshot,
    Complexity.TM.IsTransducer.reachesIn_succ_le_transducerConfigBound,
    Complexity.TM.IsTransducer.reachesIn_le_transducerConfigBound,
    Complexity.TM.DecidesInSpace.decidesInTime_configBound,
    Complexity.TM.ComputesInSpace.computesInTime_configBound}
  \leanok
  \uses{def:space-decides-in-space, def:tm-decides, def:tm-computes}
  Let $M$ be a deterministic transducer with state set $Q$ and $k$ work tapes,
  and let
  $B(n, s) = 8\,|Q|\,(n + s + 2)\,\bigl((s + 1)\,4^{s + 1}\bigr)^k$, the exact
  number of reduced snapshots (state, input head, bounded view of each work
  tape, whether the output head is on the left marker, and the symbol under the
  output head). If every configuration $M$ reaches from its initial
  configuration on $x$ is within auxiliary space $s$, and $M$ reaches a halted
  configuration in $t$ steps, then $t + 1 \le B(|x|, s)$; in particular
  $t \le B(|x|, s)$. Consequently, if $M$ decides $L$ in space $S$ then $M$
  decides $L$ in time $n \mapsto B(n, S(n))$, and if $M$ computes $f$ in space
  $S$ then $M$ computes $f$ in time $n \mapsto B(n, S(n))$.
\end{theorem}
\begin{proof}
  \leanok
  For a one-way output tape the written prefix cannot influence future
  execution, so a halting run cannot repeat a reduced snapshot.
\end{proof}

\begin{theorem}[$\cc{L} \subseteq \cc{P}$ and $\cc{FL} \subseteq \cc{FP}$]
  \label{thm:space-L-subset-P}
  \lean{Complexity.TM.transducerConfigBound_bigO_polynomial,
    Complexity.L_subset_P, Complexity.FL_subset_FP}
  \leanok
  \uses{def:L, def:FL, def:P, def:FP}
  If $S = O(\log n)$, then for every DTM with $k$ work tapes there is $c$ with
  $B(n, S(n)) = O(n^{1 + (2c + 1)k})$, where $B$ is the snapshot count of
  Theorem~\ref{thm:space-space-to-time}. Hence $\cc{L} \subseteq \cc{P}$ and
  $\cc{FL} \subseteq \cc{FP}$.
\end{theorem}
\begin{proof}
  \leanok
  \uses{thm:space-space-to-time}
  The same machine runs within the reduced-snapshot bound, which is polynomial
  for logarithmic space. The same argument places every language of $\cc{L}$
  in $\bigcup_k \cc{DTISP}(n^k, \log n) \subseteq \cc{SC}$, with the same
  machine as witness; that corollary is not yet stated in the library.
\end{proof}

\begin{theorem}[$\cc{PSPACE} \subseteq \cc{EXP}$]
  \label{thm:space-pspace-subset-exp}
  \lean{Complexity.PSPACE_subset_EXP, Complexity.TM.spaceTimeBound,
    Complexity.TM.decidesInTime_of_decidesInSpace,
    Complexity.TM.spaceTimeBound_bigO}
  \leanok
  \uses{def:PSPACE, def:EXP, def:space-decides-in-space, def:tm-decides}
  $\cc{PSPACE} \subseteq \cc{EXP}$. More precisely, if a DTM $M$ with state set
  $Q$ and $k$ work tapes decides $L$ in space $f$, transducer or not, then $M$
  decides $L$ in time
  $n \mapsto |Q|\,(n + f(n) + 2)\,\bigl((f(n) + 1)\,4^{f(n) + 1}\bigr)^k\,
  (f(n) + 2)\,4^{f(n) + 2}$, and if $f = O(n^m)$ this bound is
  $O(2^{n^j})$ for some $j$.
\end{theorem}
\begin{proof}
  \leanok
  \uses{def:space-bounded-configs}
  Every configuration of the run is windowed, so distinct configurations of
  the run have distinct codes, and a deterministic halting run visits no
  configuration twice before it halts. The halting time is therefore below the
  number of codes, and the same machine decides the language in exponential
  time.
\end{proof}

\section{Nondeterministic logarithmic space}

\begin{lemma}[NL as bounded reachability]
  \label{lem:space-nl-reachability}
  \lean{Complexity.mem_iff_exists_accepting_reachable,
    Complexity.NL_bounded_reachability,
    Complexity.NL_complement_characterization}
  \leanok
  \uses{def:NL, def:space-config-graph}
  If an NTM $M$ decides $L$ in some space bound, then $x \in L$ if and only if
  some configuration reachable from $\mathrm{init}(x)$ in the configuration
  graph of $M$ is halted with verdict $1$. For every $L \in \cc{NL}$ there are
  an NTM $M$ and constants $A, B$ such that for every $x$: $x \in L$ if and
  only if some configuration in
  $\mathrm{reachSet}(\mathrm{init}(x), A(|x| + 1)^B)$ is halted with verdict
  $1$. Equivalently, $x \notin L$ if and only if no configuration of that set is
  halted with verdict $1$.
\end{lemma}
\begin{proof}
  \leanok
  \uses{def:space-bounded-configs}
  A log-space machine has polynomially many configuration codes inside its
  window, so breadth-first search saturates within polynomially many rounds.
\end{proof}

\begin{theorem}[$\cc{NL} \subseteq \cc{P}$]
  \label{thm:space-NL-subset-P}
  \lean{Complexity.NL_subset_P, Complexity.coNL_subset_P,
    Complexity.NL_subset_P_of_search}
  \leanok
  \uses{def:NL, def:coNL, def:P, def:space-config-graph}
  $\cc{NL} \subseteq \cc{P}$ and $\cc{coNL} \subseteq \cc{P}$. The library also
  records a reduction of $\cc{NL} \subseteq \cc{P}$ to one machine, which the
  proof below does not use: if for every NTM $M$ deciding some language in
  space $S = O(\log n)$ and all constants $A, B$ some DTM decides, within an
  explicit polynomial time bound, the set of $x$ such that some configuration
  in $\mathrm{reachSet}(\mathrm{init}(x), A(|x| + 1)^B)$ of $M$ is halted with
  verdict $1$, then $\cc{NL} \subseteq \cc{P}$.
\end{theorem}
\begin{proof}
  \leanok
  \uses{lem:space-nl-reachability, thm:cobham, lem:polytime-iterate,
    lem:polytime-decision-fn, prop:polytime-P-boolean}
  The worklist search over encoded configurations is programmed in Cobham's
  function algebra and run for polynomially many rounds, and the verdict
  function gives membership in $\cc{P}$. The $\cc{coNL}$ case follows because
  $\cc{P}$ is closed under complement; it does not use the
  Immerman--Szelepcs\'enyi theorem.
\end{proof}

\begin{lemma}[Inductive-counting certificate]
  \label{lem:space-inductive-counting}
  \lean{Complexity.inductive_counting_certificate,
    Complexity.NL_membership_is_walk, Complexity.NL_nonmembership_by_counting,
    Complexity.logWindow, Complexity.NL_complement_certificate}
  \leanok
  \uses{def:space-config-graph, def:space-bounded-configs, def:NL}
  (1) If $T \subseteq \mathrm{reachSet}(c_0, i)$ and
  $|\mathrm{reachSet}(c_0, i)| \le |T|$ then $T = \mathrm{reachSet}(c_0, i)$.
  The coded search starts from one configuration code and in each round adds
  the codes of the successors of the configurations the current codes decode
  to. (2) A code lies in round $i$ of the coded search from $a_0$ if and only if
  some sequence of $i$ steps from $a_0$, each staying put or moving to such a
  successor code, reaches it. (3) If a list of distinct members of round $i$
  is at least as long as the round, every code missing from the list is
  outside the round. (4) For every $L \in \cc{NL}$ there are an NTM and
  constants $C, D, A, B$ such that $x \notin L$ if and only if round
  $A(|x| + 1)^B$ of the coded search from the code of $\mathrm{init}(x)$, in
  the window $C \log |x| + D$, can be listed in this sense with no member
  decoding to a halted configuration with verdict $1$.
\end{lemma}
\begin{proof}
  \leanok
  \uses{lem:space-nl-reachability}
  A subset of a finite set that is at least as large is the whole set; the
  codes are the fixed-width configuration codes of
  Definition~\ref{def:space-bounded-configs}.
\end{proof}

\begin{proposition}[The two inclusions are equivalent]
  \label{prop:space-nl-conl-symmetric}
  \lean{Complexity.coNL_subset_NL_of_NL_subset_coNL,
    Complexity.NL_subset_coNL_of_coNL_subset_NL,
    Complexity.NL_eq_coNL_of_NL_subset_coNL}
  \leanok
  \uses{def:NL, def:coNL}
  $\cc{NL} \subseteq \cc{coNL}$ if and only if $\cc{coNL} \subseteq \cc{NL}$.
  In particular, if $\cc{NL} \subseteq \cc{coNL}$ then $\cc{NL} = \cc{coNL}$.
\end{proposition}
\begin{proof}
  \leanok
  Complement both sides of the inclusion.
\end{proof}

\begin{proposition}[NL is closed under complement, given a counting machine]
  \label{prop:space-nl-subset-conl-of-counting}
  \lean{Complexity.NL_subset_coNL_of_counting}
  \leanok
  \uses{def:NL, def:coNL, def:space-config-graph}
  Suppose that for every NTM $M$ deciding some language in space $S$ with
  $S = O(\log n)$, and all constants $A, B$, there are constants $C, D$ and a
  nondeterministic transducer deciding, in space $C \log n + D$, the set of
  inputs $x$ such that no configuration in
  $\mathrm{reachSet}(\mathrm{init}(x), A(|x| + 1)^B)$ of $M$ is halted with
  verdict $1$. Then $\cc{NL} \subseteq \cc{coNL}$.
\end{proposition}
\begin{proof}
  \leanok
  \uses{lem:space-nl-reachability}
  Apply the hypothesis to the machine and bounds supplied by
  Lemma~\ref{lem:space-nl-reachability}; the decided set is exactly the
  complement.
\end{proof}

\begin{theorem}[Immerman--Szelepcs\'enyi]
  \label{thm:space-immerman}
  \uses{def:NL, def:coNL}
  $\cc{NL} = \cc{coNL}$.
\end{theorem}
\begin{proof}
  \uses{prop:space-nl-subset-conl-of-counting, prop:space-nl-conl-symmetric,
    lem:space-inductive-counting}
  It remains to build the counting machine required by
  Proposition~\ref{prop:space-nl-subset-conl-of-counting}: a nondeterministic
  log-space machine that guesses the certificate of
  Lemma~\ref{lem:space-inductive-counting} entry by entry, with three nested
  bounded loops (rounds, codes of a round, steps of a walk) over a constant
  number of logarithmically wide registers. The planned route builds it
  deterministically over a guess tape and converts it to an NTM. The
  statement $\cc{NL} \subseteq \cc{coNL}$ is recorded as the proposition
  \texttt{NLSubsetCoNL}.
\end{proof}

\section{Savitch's theorem and the iteration lever}

\begin{lemma}[Iterating a polynomial-time function in polynomial space]
  \label{lem:space-iterate}
  \lean{Complexity.SpaceIter.mem_PSPACE_of_iterate}
  \leanok
  \uses{def:FP, def:PSPACE, def:pair}
  Let $G \in \cc{FP}$, let $r, w$ be polynomials with natural-number
  coefficients, and let $N : \{0,1\}^* \to \mathbb{N}$ satisfy
  $1 \le N(x) \le 2^{w(|x|)}$. Write $y_i = G^i(\mathrm{pair}(\varepsilon, x))$.
  Suppose $|y_i| \le r(|x|)$ for $i \le N(x) + 1$, $y_i$ does not begin with
  $1$ (it is empty or begins with $0$) for $0 < i < N(x)$, $y_{N(x)}$ begins
  with $1$ (the done flag), and $y_{N(x)+1}$ is nonempty. If $x \in L$ exactly
  when $y_{N(x)+1}$ begins with $1$, then $L \in \cc{PSPACE}$.
\end{lemma}
\begin{proof}
  \leanok
  \uses{lem:space-poly-window}
  The machine stores only the current state and a binary iteration counter of
  polynomially many bits, applies $G$ in place until the done flag appears, and
  publishes the verdict after one more application. This is the
  polynomial-space counterpart of the iteration lemma for $\cc{FP}$ in Cobham's
  algebra.
\end{proof}

\begin{lemma}[Savitch's recursion]
  \label{lem:space-savitch-recursion}
  \lean{Complexity.savitch_halving, Complexity.savitch_reaches_within_codes,
    Complexity.NPSPACE_bounded_reachability}
  \leanok
  \uses{def:space-config-graph, def:NPSPACE}
  (1) For an NTM, $c$ reaches $c'$ within $2^{i+1}$ steps if and only if some
  midpoint is reachable from $c$ within $2^i$ steps and reaches $c'$ within
  $2^i$ steps. (2) If a map into a finite type of cardinality at most $N$
  separates the configurations reachable from $c_0$, then $c$ is reachable
  from $c_0$ if and only if it is reachable within $N$ steps. (3) For every
  $L \in \cc{NPSPACE}$ there are an NTM and a polynomial $q$ with
  natural-number coefficients such that $x \in L$ if and only if some halted
  configuration with verdict $1$ is reachable from $\mathrm{init}(x)$ within
  $2^{q(|x|)}$ steps.
\end{lemma}
\begin{proof}
  \leanok
  \uses{def:space-bounded-configs, lem:space-nl-reachability}
  A walk longer than the number of codes repeats a configuration, and the loop
  can be cut out; halving splits a walk at its midpoint. For (3), membership
  is reachability of an accepting configuration, and the configuration codes
  number at most $2$ to a polynomial.
\end{proof}

\begin{theorem}[Savitch]
  \label{thm:space-savitch}
  \lean{Complexity.NPSPACE_subset_PSPACE, Complexity.PSPACE_eq_NPSPACE}
  \leanok
  \uses{def:NPSPACE, def:PSPACE}
  $\cc{NPSPACE} \subseteq \cc{PSPACE}$, hence $\cc{PSPACE} = \cc{NPSPACE}$.
\end{theorem}
\begin{proof}
  \leanok
  \uses{lem:space-savitch-recursion, lem:space-iterate,
    lem:space-dspace-nspace, def:space-bounded-configs, thm:cobham}
  Savitch's recursion is written as a stack machine whose single step is a
  polynomial-time function, built in Cobham's algebra, on a polynomially long
  encoded stack; the midpoint enumeration ranges over every string of the
  configuration-code width. The step is iterated $2^{\mathrm{poly}}$ times by
  Lemma~\ref{lem:space-iterate}. The reverse inclusion is the embedding of
  DTMs into NTMs. No bound of the form $O(S^2)$ is extracted.
\end{proof}

\begin{theorem}[Savitch, parametric form]
  \label{thm:space-savitch-general}
  \uses{def:dspace, def:nspace, def:space-constructible}
  \notready
  If $S$ is space-constructible and $S(n) \ge \log n$, then
  $\cc{NSPACE}(S) \subseteq \cc{DSPACE}(S^2)$.
\end{theorem}
\begin{proof}
  \uses{lem:space-savitch-recursion, def:space-bounded-configs}
  The formalized theorem is only the polynomial-union corollary. The
  parametric form needs the recursion to run in space $O(S^2)$ with $S$
  computed in space $O(S)$, rather than through a polynomial-time step
  function.
\end{proof}

\begin{theorem}[Polynomial-depth game trees are evaluated in PSPACE]
  \label{thm:space-game-tree}
  \uses{def:FP, def:PSPACE}
  Fix polynomials $d$ and $\ell$. Suppose that on input $x$ a game tree has
  nodes that are strings of length at most $\ell(|x|)$ and depth at most
  $d(|x|)$; the children of a node are enumerated in order by a polynomial-time
  successor function, so the branching may be exponential; leaves carry
  polynomial-time values; and each internal node folds its children's values
  with a polynomial-time operation (such as maximum, minimum, sum, $\vee$ or
  $\wedge$) whose intermediate values have polynomial length. Then every
  language defined by a polynomial-time predicate of the root value is in
  $\cc{PSPACE}$.
\end{theorem}
\begin{proof}
  \uses{lem:space-iterate}
  Depth-first evaluation with a stack of at most $d(|x|)$ frames, each holding
  a node, a child cursor and an accumulator. One step of the walk is
  polynomial time, so Lemma~\ref{lem:space-iterate} applies. Instances:
  Savitch's theorem (an $\vee$ over midpoints of an $\wedge$ of two subcalls),
  $\cc{IP} \subseteq \cc{PSPACE}$ (maximum over prover replies, sum over
  verifier messages), $\cc{PH} \subseteq \cc{PSPACE}$ ($\vee$ and $\wedge$ over
  witnesses), $\cc{PP} \subseteq \cc{PSPACE}$ (a sum over choice sequences),
  $\cc{TQBF} \in \cc{PSPACE}$, and $\cc{NP} \cup \cc{coNP} \subseteq
  \cc{PSPACE}$, which is not yet stated. Today Savitch's theorem and
  $\cc{IP} \subseteq \cc{PSPACE}$ each program their own stack walk on
  Lemma~\ref{lem:space-iterate}, and $\cc{PH} \subseteq \cc{PSPACE}$ and
  $\cc{PP} \subseteq \cc{PSPACE}$ assemble dedicated machines; one generic
  theorem would subsume all four.
\end{proof}

\section{Quantified Boolean formulas and alternation}

\begin{definition}[Quantified Boolean formulas]
  \label{def:space-qbf}
  \lean{Complexity.QBF, Complexity.QBF.eval, Complexity.QBF.freeVars,
    Complexity.QBF.Closed, Complexity.QBF.IsTrue, Complexity.QBF.quantDepth,
    Complexity.QBF.QuantifierFree}
  \leanok
  A QBF is built from variables $x_i$ ($i \in \mathbb{N}$), the constants
  $\top$ and $\bot$, negation, conjunction, disjunction, and the quantifiers
  $\exists x_i$ and $\forall x_i$. It is evaluated under an assignment
  $\alpha : \mathbb{N} \to \{0,1\}$; a quantifier over $x_i$ evaluates its body
  under both updates of $\alpha$ at $i$ and takes the disjunction ($\exists$)
  or conjunction ($\forall$) of the two values. A formula is closed when it has
  no free variables, and true when it evaluates to $1$ under the all-false
  assignment. The quantifier depth is the largest number of quantifiers on a
  root-to-leaf path; it bounds, but does not count, alternations, and a
  formula is quantifier-free when its depth is $0$.
\end{definition}

\begin{lemma}[QBF semantics]
  \label{lem:space-qbf-semantics}
  \lean{Complexity.QBF.eval_ex_iff, Complexity.QBF.eval_all_iff,
    Complexity.QBF.eval_eq_of_agree, Complexity.QBF.eval_closed_eq,
    Complexity.QBF.eval_update_not_mem, Complexity.QBF.eval_ex_not_mem,
    Complexity.QBF.eval_all_not_mem, Complexity.QBF.eval_neg_ex,
    Complexity.QBF.eval_neg_all, Complexity.QBF.isTrue_neg_iff,
    Complexity.QBF.isTrue_conj_iff, Complexity.QBF.isTrue_disj_iff,
    Complexity.QBF.isTrue_ex_of_isTrue_all, Complexity.QBF.quantifierFree_conj}
  \leanok
  \uses{def:space-qbf}
  $\exists x_i\, \varphi$ is true under $\alpha$ if and only if $\varphi$ is
  true under $\alpha[i \mapsto b]$ for some $b$, and dually for $\forall$.
  Evaluation depends only on the free variables, so a closed formula has an
  assignment-independent value, and updating or quantifying a variable that is
  not free changes nothing. Under every assignment,
  $\neg \exists x_i\, \varphi$ has the value of $\forall x_i\, \neg \varphi$,
  and $\neg \forall x_i\, \varphi$ that of $\exists x_i\, \neg \varphi$. Truth
  commutes with negation, conjunction and disjunction, and if
  $\forall x_i\, \varphi$ is true then so is $\exists x_i\, \varphi$. A
  conjunction is quantifier-free exactly when both conjuncts are.
\end{lemma}
\begin{proof}
  \leanok
  Structural induction on the formula.
\end{proof}

\begin{definition}[TQBF]
  \label{def:space-tqbf}
  \uses{def:space-qbf, def:language}
  A bit-string codec for QBF with polynomial-time decoding and a polynomial
  size bound, and the language $\cc{TQBF}$ of codes of closed true formulas.
\end{definition}

\begin{theorem}[$\cc{TQBF} \in \cc{PSPACE}$]
  \label{thm:space-tqbf-in-pspace}
  \uses{def:space-tqbf, def:PSPACE}
  \notready
  $\cc{TQBF} \in \cc{PSPACE}$.
\end{theorem}
\begin{proof}
  \uses{thm:space-game-tree, lem:space-qbf-semantics}
  Each quantifier is a binary node of a game tree of depth at most the formula
  size.
\end{proof}

\begin{theorem}[TQBF is PSPACE-complete]
  \label{thm:space-tqbf-complete}
  \uses{def:space-tqbf, def:PSPACE, def:reduction}
  \notready
  $\cc{TQBF} \in \cc{PSPACE}$, and every $L \in \cc{PSPACE}$ satisfies
  $L \le_p \cc{TQBF}$ under polynomial-time many-one reductions.
\end{theorem}
\begin{proof}
  \uses{thm:space-tqbf-in-pspace, def:space-bounded-configs,
    lem:space-poly-window}
  Hardness: express ``$c$ reaches $c'$ within $2^i$ steps'' by a formula that
  quantifies universally over the pair of halves, so its size is linear in $i$
  rather than exponential, over a formula for one step between configuration
  codes.
\end{proof}

\begin{definition}[Alternating machines]
  \label{def:space-atm}
  \uses{def:tm, def:ntm}
  An alternating TM partitions its non-halting states into existential and
  universal states; acceptance is the least fixed point of the induced
  $\vee$/$\wedge$ labelling of the configuration graph, or equivalently a
  bounded game between two players. $\cc{AP}$ is the class of languages decided
  by an alternating machine all of whose paths halt in polynomial time.
\end{definition}

\begin{theorem}[$\cc{AP} = \cc{PSPACE}$]
  \label{thm:space-ap-eq-pspace}
  \uses{def:space-atm, def:PSPACE}
  \notready
  $\cc{AP} = \cc{PSPACE}$.
\end{theorem}
\begin{proof}
  \uses{thm:space-game-tree, thm:space-tqbf-complete}
  $\cc{AP} \subseteq \cc{PSPACE}$ is an instance of game-tree evaluation;
  $\cc{PSPACE} \subseteq \cc{AP}$ follows by evaluating TQBF alternately, or
  directly by alternating Savitch's midpoint recursion.
\end{proof}

\section{Space hierarchy}

\begin{definition}[Space-constructible functions]
  \label{def:space-constructible}
  \uses{def:tm}
  $S$ is space-constructible if some DTM maps $1^n$ to the binary
  representation of $S(n)$ in space $O(S(n))$. This is distinct from the
  library's time-constructibility and from the randomness-bound condition used
  in the PCP theorem.
\end{definition}

\begin{theorem}[Deterministic space hierarchy]
  \label{thm:space-hierarchy-det}
  \uses{def:dspace, def:space-constructible}
  \notready
  If $S'$ is space-constructible, $S'(n) \ge \log n$ and $S = o(S')$, then
  $\cc{DSPACE}(S) \subsetneq \cc{DSPACE}(S')$.
\end{theorem}
\begin{proof}
  \uses{thm:time-hierarchy, thm:space-space-to-time}
  Diagonalize against machines with a space-bounded universal simulation,
  reusing the diagonal-machine infrastructure of the time hierarchy theorem.
\end{proof}

\begin{theorem}[Nondeterministic space hierarchy]
  \label{thm:space-hierarchy-nondet}
  \uses{def:nspace, def:space-constructible}
  \notready
  If $S'$ is space-constructible, $S'(n) \ge \log n$ and $S = o(S')$, then
  $\cc{NSPACE}(S) \subsetneq \cc{NSPACE}(S')$.
\end{theorem}
\begin{proof}
  \uses{thm:space-immerman, thm:space-hierarchy-det}
  Closure of nondeterministic space under complement, by inductive counting,
  makes the diagonal argument go through.
\end{proof}

\section{A log-space programming layer}

Membership in $\cc{FL}$ is currently established machine by machine, for
example for the tableau serializer behind logspace-uniform circuits, and the
only structural fact about $\cc{FL}$ is $\cc{FL} \subseteq \cc{FP}$. Two
layers already support writing such machines. Space-aware Hoare contracts
(Definition~\ref{def:space-hoare}) pair a time-bounded triple with an
all-reachable auxiliary-space bound, compose sequentially at the larger of the
two budgets, and package per-input contracts of a transducer into a
$\mathrm{ComputesInSpace}$ statement. An experimental layer of proof-carrying
binary register routines (\texttt{BinaryRoutine}) composes sequentially,
supports loops with width certificates that collapse to logarithmic space, and
yields a $\mathrm{ComputesInSpace}$ contract; its programs see only the input
length, not the input bits. The planned nodes below would make log-space
programming as routine as Cobham's algebra makes polynomial time.

\begin{definition}[Space-aware Hoare contracts]
  \label{def:space-hoare}
  \lean{Complexity.TM.HoareSpace, Complexity.TM.HoareTimeSpace}
  \leanok
  \uses{def:models-hoare, def:space-decides-in-space}
  For a DTM $M$, a precondition $\mathit{pre}$ on (input tape, work tapes,
  output tape), an input length $n$ and a space bound $s$, the space contract
  holds if, whenever $M$ starts in its start state on tapes satisfying
  $\mathit{pre}$, every configuration it reaches is within auxiliary space $s$
  for input length $n$. The time-and-space contract with bounds $b$, $n$ and
  $s$ is the conjunction of the time-bounded triple
  $\{\mathit{pre}\}\, M\, \{\mathit{post}\}_b$ with this space contract.
\end{definition}

\begin{lemma}[Rules for space-aware contracts]
  \label{lem:space-hoare-rules}
  \lean{Complexity.TM.hoareSpace_of_invariant,
    Complexity.TM.HoareTime.toHoareTimeSpace,
    Complexity.TM.HoareTimeSpace.consequence, Complexity.TM.IsTransducer.seqTM,
    Complexity.TM.seqTM_hoareTimeSpace,
    Complexity.TM.computesInSpace_of_hoareTimeSpace}
  \leanok
  \uses{def:space-hoare, def:models-combinators, def:space-decides-in-space}
  (1) A configuration predicate that holds at every start satisfying
  $\mathit{pre}$, is preserved by every step, and implies auxiliary space $s$
  gives the space contract. (2) If $\{\mathit{pre}\}\, M\,
  \{\mathit{post}\}_b$ holds and every start configuration satisfying
  $\mathit{pre}$ is within auxiliary space $s_0$, then the time-and-space
  contract holds with space $s_0 + b$. (3) A time-and-space contract survives
  strengthening the precondition, weakening the postcondition, and enlarging
  the time bound, the input length and the space bound. (4) The sequential
  composition of two transducers is a transducer. If $M_1$ satisfies the
  time-and-space contract from $\mathit{pre}$ to $\mathit{mid}$ with bounds
  $b_1, n, s_1$, the phase-boundary normalization maps tapes satisfying
  $\mathit{mid}$ to tapes satisfying $\mathit{mid}'$, and $M_2$ satisfies the
  contract from $\mathit{mid}'$ to $\mathit{post}$ with bounds $b_2, n, s_2$,
  then $\mathrm{seq}(M_1, M_2)$ satisfies the contract from $\mathit{pre}$ to
  $\mathit{post}$ with bounds $b_1 + 1 + b_2$, $n$ and $\max(s_1, s_2)$.
  (5) If a transducer satisfies, for every input $x$, the time-and-space
  contract from the initial tapes on $x$ to ``the output tape holds $f(x)$''
  with bounds $T(|x|)$, $|x|$ and $S(|x|)$, then it computes $f$ in space $S$.
\end{lemma}
\begin{proof}
  \leanok
  \uses{thm:models-seq-hoare}
  (1) is induction along the reachable configurations. For (2), the run halts
  within $b$ steps, so every reachable configuration is at most $b$ steps from
  its start, and each step moves each head by at most one cell. For (4), the
  time component is
  Theorem~\ref{thm:models-seq-hoare}; the first phase is covered by the
  contract of $M_1$, and the normalized start of $M_2$ and everything after it
  by the contract of $M_2$.
\end{proof}

\begin{theorem}[FL is closed under composition]
  \label{thm:space-fl-comp}
  \uses{def:FL, def:L}
  If $f, g \in \cc{FL}$ then $g \circ f \in \cc{FL}$. If moreover
  $A \in \cc{L}$, then $f^{-1}(A) \in \cc{L}$.
\end{theorem}
\begin{proof}
  \uses{thm:space-L-subset-P}
  The recomputation argument: simulate the machine for $g$, keeping the
  position of its virtual input head in a binary counter, and whenever it reads
  bit $j$ of $f(x)$, rerun the machine for $f$ from the start, discarding output
  until position $j$. By $\cc{FL} \subseteq \cc{FP}$, $|f(x)|$ is polynomial, so
  the counters need $O(\log n)$ bits. Under the library's convention the
  virtual input head may travel one cell past $f(x)$ for free and is charged
  beyond that.
\end{proof}

\begin{proposition}[Log-space reductions]
  \label{prop:space-logspace-reductions}
  \uses{def:FL, def:L, def:NL, def:reduction}
  Write $A \le_{\log} B$ when some $f \in \cc{FL}$ satisfies
  $x \in A \iff f(x) \in B$. Then $\le_{\log}$ is reflexive and transitive,
  $A \le_{\log} B$ implies $A \le_p B$, and $\cc{L}$ and $\cc{NL}$ are closed
  downward under $\le_{\log}$.
\end{proposition}
\begin{proof}
  \uses{thm:space-fl-comp, thm:space-L-subset-P}
  Transitivity and closure of $\cc{L}$ are Theorem~\ref{thm:space-fl-comp};
  closure of $\cc{NL}$ is the same recomputation with a nondeterministic outer
  machine.
\end{proof}

\begin{theorem}[PATH is NL-complete]
  \label{thm:space-path-nl-complete}
  \uses{prop:space-logspace-reductions, def:NL}
  \notready
  Directed $s$--$t$ reachability, under an adjacency-matrix codec, is in
  $\cc{NL}$ and every language in $\cc{NL}$ log-space reduces to it.
\end{theorem}
\begin{proof}
  \uses{def:space-config-graph, lem:space-nl-reachability}
  The reduction writes out the configuration graph of the log-space machine on
  the given input, one adjacency bit at a time.
\end{proof}

\begin{definition}[Log-space loop programs]
  \label{def:space-logspace-programs}
  \uses{def:FL}
  A program has a fixed finite set of registers, each holding a natural number
  of at most $c \log n$ bits for a program constant $c$. Instructions are
  read-only access to the input bit at a register-held index, the input length,
  successor, comparison and addition on registers, conditionals, bounded loops
  whose bounds are register values, sequencing, and appending a bit to a
  write-only output stream.
\end{definition}

\begin{theorem}[Compiling loop programs to log-space transducers]
  \label{thm:space-logspace-compiler}
  \uses{def:space-logspace-programs, def:FL, def:L}
  \notready
  The function computed by every log-space loop program is in $\cc{FL}$, and a
  program that outputs a single bit decides a language in $\cc{L}$.
\end{theorem}
\begin{proof}
  \uses{thm:space-fl-comp, lem:space-hoare-rules}
  Each register becomes a binary work tape, input access becomes a seek driven
  by a counter, and loops become binary counters, reusing the Hoare-style
  binary-counter subroutines and the space-aware contracts of
  Lemma~\ref{lem:space-hoare-rules}. The counting machine of
  Theorem~\ref{thm:space-immerman} and the tableau serializer would then be
  written as programs.
\end{proof}
